% design_software.tex — Documento di Design del Software

\section{Design del Software}

\subsection{Schema per la persistenza dati}

Il database relazionale è gestito tramite JPA/Hibernate con strategia
\texttt{ddl-auto=update}. Di seguito lo schema entità-relazione con tutte le
tabelle, colonne e vincoli.

\subsubsection{Tabella \texttt{users}}

\begin{center}
\small
\begin{tabular}{|l|l|l|p{4cm}|}
\hline
\textbf{Colonna} & \textbf{Tipo} & \textbf{Vincoli} & \textbf{Note} \\
\hline
\texttt{id} & UUID & PK, AUTO & Identificativo univoco \\
\hline
\texttt{email} & VARCHAR & UNIQUE, NOT NULL & Indirizzo email \\
\hline
\texttt{password\_hash} & CHAR(60) & NOT NULL & Hash BCrypt \\
\hline
\texttt{name} & VARCHAR & NOT NULL & Nome visualizzato \\
\hline
\texttt{role} & VARCHAR & NOT NULL & Enum: ADMIN, USER, READONLY \\
\hline
\texttt{created\_at} & TIMESTAMP & & Data di registrazione \\
\hline
\end{tabular}
\end{center}

\subsubsection{Tabella \texttt{bugs}}

\begin{center}
\small
\begin{tabular}{|l|l|l|p{4cm}|}
\hline
\textbf{Colonna} & \textbf{Tipo} & \textbf{Vincoli} & \textbf{Note} \\
\hline
\texttt{id} & UUID & PK, AUTO & Identificativo univoco \\
\hline
\texttt{title} & VARCHAR & NOT NULL & Titolo del bug \\
\hline
\texttt{description} & TEXT & NOT NULL & Descrizione dettagliata \\
\hline
\texttt{type} & VARCHAR & NOT NULL & Enum: QUESTION, BUG, DOCUMENTATION, FEATURE \\
\hline
\texttt{status} & VARCHAR & & Enum: TODO, IN\_PROGRESS, RESOLVED, ARCHIVED \\
\hline
\texttt{priority} & VARCHAR & & Enum: LOW, MEDIUM, HIGH, URGENT \\
\hline
\texttt{archived} & BOOLEAN & default false & Flag di archiviazione \\
\hline
\texttt{created\_at} & TIMESTAMP & & Data di creazione \\
\hline
\texttt{deadline} & TIMESTAMP & nullable & Scadenza opzionale \\
\hline
\texttt{duplicate\_of} & UUID & FK $\rightarrow$ bugs.id & Bug duplicato (self-ref) \\
\hline
\texttt{created\_by} & UUID & FK $\rightarrow$ users.id & Creatore \\
\hline
\texttt{assignee\_id} & UUID & FK $\rightarrow$ users.id & Assegnatario \\
\hline
\end{tabular}
\end{center}

\subsubsection{Tabella \texttt{comments}}

\begin{center}
\small
\begin{tabular}{|l|l|l|p{4cm}|}
\hline
\textbf{Colonna} & \textbf{Tipo} & \textbf{Vincoli} & \textbf{Note} \\
\hline
\texttt{id} & UUID & PK, AUTO & Identificativo univoco \\
\hline
\texttt{text} & TEXT & NOT NULL & Testo del commento \\
\hline
\texttt{bug\_id} & UUID & FK $\rightarrow$ bugs.id, NOT NULL & Bug associato \\
\hline
\texttt{author\_id} & UUID & FK $\rightarrow$ users.id, NOT NULL & Autore \\
\hline
\texttt{created\_at} & TIMESTAMP & & Data di creazione \\
\hline
\end{tabular}
\end{center}

\subsubsection{Tabella \texttt{history}}

\begin{center}
\small
\begin{tabular}{|l|l|l|p{4cm}|}
\hline
\textbf{Colonna} & \textbf{Tipo} & \textbf{Vincoli} & \textbf{Note} \\
\hline
\texttt{id} & UUID & PK, AUTO & Identificativo univoco \\
\hline
\texttt{bug\_id} & UUID & FK $\rightarrow$ bugs.id, NOT NULL & Bug associato \\
\hline
\texttt{who\_id} & UUID & FK $\rightarrow$ users.id, NOT NULL & Utente esecutore \\
\hline
\texttt{action} & VARCHAR & NOT NULL & Enum: CREATE, UPDATE, ASSIGN, COMMENT, ARCHIVE, DUPLICATE, LABELS\_SET, ATTACHMENT\_UPLOADED \\
\hline
\texttt{details} & TEXT & nullable & Dettagli del cambiamento \\
\hline
\texttt{created\_at} & TIMESTAMP & & Timestamp dell'azione \\
\hline
\end{tabular}
\end{center}

\subsubsection{Tabella \texttt{labels}}

\begin{center}
\small
\begin{tabular}{|l|l|l|p{4cm}|}
\hline
\textbf{Colonna} & \textbf{Tipo} & \textbf{Vincoli} & \textbf{Note} \\
\hline
\texttt{id} & UUID & PK, AUTO & Identificativo univoco \\
\hline
\texttt{name} & VARCHAR & UNIQUE, NOT NULL & Nome della label \\
\hline
\texttt{created\_at} & TIMESTAMP & & Data di creazione \\
\hline
\end{tabular}
\end{center}

\subsubsection{Tabella di join \texttt{bug\_labels}}

\begin{center}
\small
\begin{tabular}{|l|l|l|p{5cm}|}
\hline
\textbf{Colonna} & \textbf{Tipo} & \textbf{Vincoli} & \textbf{Note} \\
\hline
\texttt{bug\_id} & UUID & FK $\rightarrow$ bugs.id & Relazione molti-a-molti Bug $\leftrightarrow$ Label \\
\hline
\texttt{label\_id} & UUID & FK $\rightarrow$ labels.id & \\
\hline
\end{tabular}
\end{center}

\subsubsection{Diagramma E-R semplificato}

\begin{verbatim}
  users 1 ---< N bugs        (createdBy)
  users 1 ---< N bugs        (assignee)
  bugs  1 ---< N comments
  bugs  1 ---< N history
  users 1 ---< N comments    (author)
  users 1 ---< N history     (who)
  bugs  N >---< M labels     (bug_labels)
  bugs  N --->  1 bugs       (duplicateOf, self-referential)
\end{verbatim}

% ────────────────────────────────────────────────────────

\subsection{Diagramma delle classi di design}

Il diagramma delle classi di design è organizzato secondo il pattern architetturale
\textbf{layered architecture} con i seguenti package:

\begin{description}
    \item[\texttt{model}] Entità JPA che mappano le tabelle del database:
    \texttt{User}, \texttt{Bug}, \texttt{Comment}, \texttt{History}, \texttt{Label}
    e le enumerazioni \texttt{Role}, \texttt{BugType}, \texttt{BugStatus},
    \texttt{Priority}, \texttt{HistoryAction}.

    \item[\texttt{repository}] Interfacce Spring Data JPA che estendono
    \texttt{JpaRepository<T, UUID>}:
    \texttt{UserRepository}, \texttt{BugRepository} (+ \texttt{JpaSpecificationExecutor}),
    \texttt{CommentRepository}, \texttt{LabelRepository}, \texttt{HistoryRepository}.

    \item[\texttt{service}] Classi di business logic:
    \texttt{BugService} (gestione completa del ciclo di vita dei bug),
    \texttt{AuthService} (registrazione, login, generazione token).

    \item[\texttt{dto}] Record Java per i dati di trasferimento:
    \texttt{BugCreateRequest}, \texttt{BugPatchRequest}, \texttt{AssignRequest},
    \texttt{LabelSetRequest}, \texttt{LoginRequest}.

    \item[\texttt{controller}] Controller REST:
    \texttt{BugController} (\texttt{/api/bugs}),
    \texttt{AuthController} (\texttt{/api/auth}),
    \texttt{AdminController} (\texttt{/api/admin}),
    \texttt{ExportController} (\texttt{/api/export}).

    \item[\texttt{security}] Infrastruttura di sicurezza:
    \texttt{JwtService} (generazione/validazione token),
    \texttt{JwtAuthFilter} (filtro per estrarre e verificare il JWT).

    \item[\texttt{config}] Configurazione dell'applicazione:
    \texttt{SecurityConfig} (CORS, CSRF, filter chain),
    \texttt{DataSeeder} (dati iniziali per lo sviluppo).
\end{description}

\begin{verbatim}
  +--------------------+     +-------------------+     +------------------+
  |    Controller      |---->|     Service       |---->|    Repository    |
  |  (REST Endpoints)  |     | (Business Logic)  |     |  (Data Access)  |
  +--------------------+     +-------------------+     +------------------+
         |                           |                         |
         v                           v                         v
  +--------------------+     +-------------------+     +------------------+
  |       DTO          |     |      Model        |     |     Database     |
  | (Request/Response) |     |  (JPA Entities)   |     |  (H2/PostgreSQL) |
  +--------------------+     +-------------------+     +------------------+
                                     ^
                                     |
                             +-------------------+
                             |    Security       |
                             | (JWT, Auth Filter)|
                             +-------------------+
\end{verbatim}

% ────────────────────────────────────────────────────────

\subsection{Scelte di software design}

\subsubsection{Pattern Repository}
Ogni entità dispone di un'interfaccia repository dedicata che estende
\texttt{JpaRepository}. Questo pattern separa la logica di accesso ai dati
dalla logica di business, favorendo la testabilità tramite mock objects.

\texttt{BugRepository} estende anche \texttt{JpaSpecificationExecutor},
consentendo la costruzione di query dinamiche (filtri per tipo, stato,
priorità, label, ricerca testuale) senza codice SQL custom.

\subsubsection{Pattern Service Layer}
La logica di business è concentrata nei servizi (\texttt{BugService},
\texttt{AuthService}), che orchestrano le operazioni tra più repository.
I controller sono \emph{thin}: delegano immediatamente ai servizi senza
contenere logica di business.

\subsubsection{Record Java per i DTO}
I DTO sono implementati come \texttt{record} Java, che garantiscono
immutabilità, auto-generazione di \texttt{equals}/\texttt{hashCode}/\texttt{toString},
e riduzione del codice boilerplate. Le annotazioni di validazione
(\texttt{@NotBlank}, \texttt{@NotNull}, \texttt{@NotEmpty}) sono applicate
direttamente sui componenti del record.

\subsubsection{Specification Pattern per le query}
La classe \texttt{BugSpecification} implementa il pattern Specification
di Spring Data, componendo predicati JPA Criteria in modo dichiarativo.
I filtri (tipo, stato, priorità, label, ricerca full-text su titolo e
descrizione) possono essere combinati liberamente dal frontend.

\subsubsection{Audit Trail tramite History}
Ogni operazione significativa (creazione, modifica, assegnazione,
commento, archiviazione) genera automaticamente una entry nella tabella
\texttt{history}, con riferimento all'utente, all'azione e ai dettagli
della modifica. Questo garantisce tracciabilità completa senza librerie esterne.

\subsubsection{Autenticazione stateless con JWT}
\texttt{JwtService} genera token firmati HMAC-SHA256 con claims personalizzati
(ruolo, nome, email). \texttt{JwtAuthFilter} intercetta ogni richiesta,
estrae il token da cookie o header, e popola il \texttt{SecurityContext}
di Spring Security. Questo approccio elimina la necessità di sessioni
server-side e facilita la scalabilità orizzontale.

\subsubsection{Gestione delle autorizzazioni}
Il controllo degli accessi è implementato a due livelli:
\begin{itemize}
    \item \textbf{Livello HTTP}: Spring Security filtra le richieste in base
    all'autenticazione (JWT valido).
    \item \textbf{Livello Service}: le regole di business (solo admin può
    assegnare; solo assegnatario o admin può cambiare stato) sono verificate
    nel servizio, lanciando \texttt{SecurityException} in caso di violazione.
\end{itemize}

% ────────────────────────────────────────────────────────

\subsection{Evidenza dell'uso di strumenti di versioning}

Il codice sorgente del progetto è ospitato su un repository \textbf{GitHub}
accessibile ai membri del team di sviluppo.

\begin{itemize}
    \item \textbf{Repository}: \url{https://github.com/alessandroilprogrammatore/BugBoard26}
    \item \textbf{Release}: la release ``Specifica dei requisiti'' è stata
    pubblicata il 19 ottobre 2025
    \item \textbf{Branching strategy}: sviluppo su branch feature,
    merge su \texttt{main} tramite pull request
    \item \textbf{Commit}: messaggi descrittivi che seguono il formato
    convenzionale (\texttt{feat:}, \texttt{fix:}, \texttt{docs:}, \texttt{test:})
\end{itemize}
